% !TEX root = ../../hw07.tex
% Source: author's supplied text in chat, 2026-11-06, Q4 (i)--(ii).
% Keep both L'Hopital calculations. Author approved the amendments, 2026-11-06.
\begin{proof}[Solution]
  Fix $x$ and take $h$ sufficiently small that
  $x-h$ and $x+h$ lie in the domain. We differentiate with respect
  to $h$, keeping $x$ fixed. Twice differentiability makes $f$ and
  $f'$ continuous, and $f''$ is continuous by assumption. We apply
  L'H\^opital's rule on each side of $h=0$.

  For part (i), since the limit is in $0/0$ form, we can apply
  L'H\^opital's rule,
  \[
    \begin{aligned}
      \lim_{h\to0}\frac{f(x+h)-f(x-h)}{2h}
        &\overset{\mathrm{L'H}}{=}
          \lim_{h\to0}\frac{f'(x+h)+f'(x-h)}2\\[4pt]
        &=\frac{f'(x)+f'(x)}2\\
        &=f'(x).
    \end{aligned}
  \]
  Here the denominator derivative is $2\ne0$,
  and continuity of $f'$ gives the derivative-ratio limit.

  Similarly, part (ii) is also in $0/0$ form,
  \[
    \begin{aligned}
      \lim_{h\to0}\frac{f(x+h)+f(x-h)-2f(x)}{h^2}
        &\overset{\mathrm{L'H}}{=}
          \lim_{h\to0}\frac{f'(x+h)-f'(x-h)}{2h}\\[4pt]
        &\overset{\mathrm{L'H}}{=}
          \lim_{h\to0}\frac{f''(x+h)+f''(x-h)}2\\[4pt]
        &=\frac{f''(x)+f''(x)}2\\
        &=f''(x).
    \end{aligned}
  \]
  For the first application, the denominator
  derivative $2h$ is nonzero when $h\ne0$. The intermediate
  quotient is again of $0/0$ form by continuity of $f'$. Its
  denominator derivative is $2$, and continuity of $f''$ supplies
  the final limit.
\end{proof}
